\begin{abstract}
Reliable maritime oil spill surveillance is fundamentally challenged by the intensity mimicry of biogenic ``lookalikes,'' such as algal blooms and low-wind calm zones. Current intensity-based segmentation models suffer from persistent false-alarm rates, leading to prohibitively expensive $O(10^5)$ USD operational response costs for non-hazardous slicks. This work integrates AQUA-SENTINEL V9, a Physics-Guided Dual-Encoder (PGDE) architecture that resolves this ambiguity by combining second-order Laplacian surface roughness gradients with multimodal neural streams. By gating the segmentation bottleneck through an active auxiliary suppression head, the framework prioritizes deployment reliability over raw baseline recall. Validated on a multimodal dataset of 7,528 samples, the PGDE achieves a peak precision of 99.66\% and an operational Dice score of 0.7922. These results provide a robust, deployment-ready methodology for autonomous environmental monitoring, effectively eliminating the geographical and economic burden of manual spill verification in global oceanic corridors.
\end{abstract}
